Tot tien trekken
Aftrekspanning
Het uitmelkmannetje

Dus de brieven bleek  tante mee te nemen voordat pa ze verbrand en tijdens de vlucht van Emmy is zij een tijdje bij tante. En tante geeft haar, als Emmy dan weer verder gaat, de brieven mee. Dat doet ze pas bij haar vlucht, omdat ze zegt, ja ik heb er spijt van dat ik dat nu pas geef. “Maar ik weet niet of de inhoud zo goed voor je is, maar ik kan hier natuurlijk niet het contact met je zus onthouden." Dus leest Emmy pas de brieven als ze op pad gaat.

de brieven worden uiteindelijk door Tannen. Dat is een monnik die in een roedel leeft waar Emmie een tijdje onderdak is. Gezien en gekleemd en ook vernietigd. Al schijnde godslasterend en niet daglicht verdragen. Emmie is het daar op zich mee eens. Maar ze heeft de brieven ondertussen in haar periode van vogelvrijheid ook meer als een soort afleiding. Want ja, er was niet veel te doen. Zeker niet in de periode dat ze alleen zwierf. Heeft ze de brieven zo vaak gelezen dat ze zich de inhoud heel goed kan herinneren in haar hoofd.

Door de inhoud van de brieven weet tannen dat ze tot een gewillige hoer kan uitgroeien

Ik breng rust in het veld grap

De monnik gebruikt de brieven als een soort afpers middel van kijk eens waar jij vandaan komt en waarom doe je nou zo moeilijk als wij je vragen om een gunst te doen voor de hoofdrolspeler Einfach. Jij hebt de genen van een hoer.

Ze moest terugdenken

aan haar medelijden
voor Einfach zijn lijden.
Die hield zo van meiden.

Dat was onomstoten
voor hem uitgesloten.
Nooit had hij genoten

van een vrouw haar aandacht.
Door haar aantrekkingskracht
ging hij uit alle macht

Emmy steeds ompraten
om zich in te laten
met zijn lustdictaten.

Het ging voor onze ster
meer dan een brug te ver.
Elk voorstel was vulgair

en daarnaast verboden
volgens haar geboden.
Zijn wervingsmethoden

ontweek Emmy minzaam.
Rustig en bedachtzaam
werden handen langzaam

weer teruggegeven.
Met veel medeleven
kon ze hem vergeven

na zichzelf vervuilen
en mocht hij zelfs schuilen
als hij dan moest huilen

in haar veilige schoot.
Waar hij zo van genoot
dat hij zich weer aanbood.

Een kansloze poging
waar zij kalm op inging
door weer een weigering

om hem te gedogen.
Met zijn lede ogen
keek uit mededogen

de Herder dit steeds aan.
Hij sprak haar erop aan.
“Laat mijn broer toch begaan”.

“Mijn geloof geeft mij raad
en verbiedt zulk een daad.
Ik snap; hij heeft het kwaad.

Ik begrijp je wens goed.
Begrijp ook mijn gemoed.
Hij weet niet wat hij doet.

Hij kan er niets aan doen,
maar de misvormde oen
kent nauwelijks fatsoen.

Al die grove zinnen
over mij beminnen
en mij overwinnen

kan ik accepteren.
Ik zal me verweren
als die impliceren

dat ik erin mee ga
en er voor open sta.
Weet dat ik ze af sla.

Ik moet me verzetten
om hem te beletten
me pardoes te pletten.

Hoewel hij van mij houdt,
zijn we ook niet getrouwd.
En hoe hij is gebouwd,

is voor mij afstotend.
Dat is niet liefkozend
en misschien zelfs stotend,

maar wel mijn beleving.
Ik voel zelfs huivering
bij een toenadering.”

De Herder zag profijt
in Emmy haar vroomheid
en een mogelijkheid

om haar toch te lijmen.
Hij begon met slijmen
door godsvrucht te rijmen

met een bravourestuk
in plaats van met groepsdruk
en angst voor ongeluk.

Ook had de slimmerik
een ijzersterke klik
met de roedels monnik.

De monnik was gevlucht
toen men in zijn gehucht
lucht kreeg van zijn ontucht

met een dienstbare non,
die hij slinks overwon.
De schuine monnik kon

de herder adoreren
en de troep bezweren
zich niet te verweren

tegen zijn dwaasheden.
Hij kon overreden
zonder goede reden.

Hij werd alom gevreesd
dat hij kalm en bedeesd
in iedereen zijn geest

wist binnen te dringen
zonder je te dwingen.
Mensen die weggingen

van hun vertrouwde plek,
die kreeg je ook niet gek
met dwang of getouwtrek.

“Emmy, mijn lieflijk kind.
De heer weet wat je vind
en is je goedgezind,

want je bent vroom en kuis.
Ik ben maar een spreekbuis
en zelf geheel niet thuis

in zaken van het vlees.
Het is God die ik vrees
en wiens boek ik na lees.

Het ‘heb uw naasten lief’
geeft mij zoveel gerief.
Is het niet een vrijbrief?

God vergeeft toch het meest
de zwakkeren van geest
is wat men apert leest?”

Talrijke citaten,
Bijbelse dictaten
en mogelijke baten

passeerden de revue.
Het punt was continu:
haar afwijzing was cru

en niet noodzakelijk.
“God neemt het niet kwalijk
als je beminnelijk

een medemens bedient
als die dat zo verdiend.
Iedere mensenvriend

krijgt een speciaal plekje
in God zijn optrekje.
Zelfs die met een vlekje…

Een man die binnendringt,
is waar de schoen pas wringt
en zeker als ie zingt

voordat zijn buis vol zaad
op tijd de kerk uitgaat.
Geldt dit voor ons gelaat,

met de mond als flacon?
De vrouw van Solomon
toonde hoe het wel kon

om genot te geven
zonder fout te leven.
Nergens staat geschreven

dat het zo ook niet mag.”
En zo sprak het gezag
tegen haar dag na dag.

De kans die Emmy bood
aan de roedeldespoot
was haar gewetensnood.

Snode woorden vol pracht
kregen Emmy’s aandacht.
De monnik kreeg de macht

over haar maakbaar brein.
Hij vermurwde haar zijn
en hakte het toen fijn.

“Je weet natuurlijk wel;
God gaf in Samuel
ons als mens het bevel

niet naar het uiterlijk,
maar naar het innerlijk
van een menselijke

creatuur te kijken.
God zal ons aanreiken
wie ons zal verrijken.”
